%% kapa: exam=1 round=33 year=2022 period=1 form=A subject=real_estate subject_ko=부동산학원론 seq=2 number=82
%% source: Q-Net 감정평가사 기출문제 (cst00309) artlSeq=5209997 file=raw/1008/5209997/2022년도 제33회 감정평가사 1차시험 1교시(민법,경제학,부동산학).pdf
%% conversion: kordoc 4.24.51 -> tools/gen1.py (auto); review=proofread-sonnet; flags=none
\begin{question}{82}
부동산활동에 관련된 설명으로 옳은 것을 모두 고른 것은?
\begin{framedtext}
ㄱ. 공유지(共有地)란 1필지의 토지를 2인 이상이 공동으로 소유한 토지로, 지분비율 또는 지분의 위치에 따라 감정평가한다. \par
ㄴ. 일단지란 용도상 불가분의 관계에 있고 지가형성요인이 같은 2필지 이상의 토지로, 필지별로 감정평가한다. \par
ㄷ. 선하지란 고압선 아래의 토지로, 고압선등 통과부분의 면적 등 제한의 정도를 고려하여 감정평가한다. \par
ㄹ. 맹지란 도로와 접한 면이 없는 토지로, 도로로 사용하기 위한 지역권이 설정되어 있는 경우 도로가 있는 것으로 보고 감정평가한다. \par
ㅁ. 환지란 도시개발사업에서 사업전 토지의 위치 등을 고려하여 소유자에게 재분배하는 사업후의 토지로, 환지처분 이전에 환지예정지로 지정된 경우에는 종전 토지의 위치 등을 기준으로 감정평가한다.
\end{framedtext}
\begin{choices}
  \item ㄱ, ㄴ, ㄷ
  \item ㄱ, ㄷ, ㄹ
  \item ㄱ, ㄷ, ㅁ
  \item ㄴ, ㄷ, ㄹ
  \item ㄴ, ㄹ, ㅁ
\end{choices}
\end{question}
\answer{2}
